% Original Markdown SHA-256: 81a329e64b8f54838228faf23ef8f9f60c2f11c9ac0d4a4067bc89b55122471b
\chapter{An unweighted ternary-variance refinement of the k = 4 lower bound}
\label{chapter:01_ternary-variance-lower-bound}
{\small\noindent\textbf{Source:} \nolinkurl{history/EP817_PR2_variance_contribution_20260914/notes/ternary-variance-lower-bound.md}\
\textbf{Identity:} \nolinkurl{81a329e64b8f54838228faf23ef8f9f60c2f11c9ac0d4a4067bc89b55122471b}}
\medskip
Date: 14 September 2026. Base: \nolinkurl{23c0110c95b5a2036bdc04a1f352b6e5e27742c8}.

\textbf{Status:} ordinary mathematical derivation, with exact bounded computational checks; submitted for independent review. \textbf{Not Lean-checked.} No existing Lean source, receipt, or accepted claim status is changed by this contribution.

The dependency is the minimal signed-block/short-kernel classification proved in \texttt{paper/ep817\_k4\_rate.tex}, especially \texttt{lem:short-kernel} and \texttt{eq:ternary-count}. That construction and proof remain attributed to the anonymous/deleted Reddit contributor. The finite upper-bound Lean extension remains separately credited to \texttt{sneed-and-feed}. The calculations below are an additional workbench contribution, not a priority claim for either earlier result. No identification of the anonymous contributor is attempted.

\section{1. Statement}\label{n01-statement}

Let A = \{a\_1,...,a\_n\} be a set of distinct positive integers, n \textgreater= 1, with no nonconstant four-term arithmetic progression in its set H(A) of subset sums. Write N = max A, S = sum\_i a\_i, Q = sum\_i a\_i\^{}2, and

\[
\begin{adjustbox}{max width=.92\linewidth}
$\displaystyle T(A)=\left\{\sum_i x_i a_i:x_i\in\{0,1,2\}\right\},\qquad M=|T(A)|.$
\end{adjustbox}
\]

Crucially, every element of T(A) is counted \textbf{once}, regardless of the number of ternary representations. Put n = 3q+r, 0 \textless= r \textless{} 3, and M\_n = 19\^{}q 3\^{}r.

\textbf{Theorem.} Under these hypotheses,

\[
\begin{adjustbox}{max width=.92\linewidth}
$\displaystyle M\ge M_n,\qquad
19(M^2-1)\le192Q\le192nN^2.$
\end{adjustbox}
\]

Consequently the extremal function in Problem 817 satisfies

\[
\begin{adjustbox}{max width=.92\linewidth}
$\displaystyle \boxed{g_4(n)\ge
\left\lceil\sqrt{\frac{19}{192n}\bigl(M_n^2-1\bigr)}\right\rceil.}$
\end{adjustbox}
\]

In particular,

\[
\begin{adjustbox}{max width=.92\linewidth}
$\displaystyle g_4(n)\ge\sqrt{\frac{19}{192n}\bigl(19^{2n/3}-1\bigr)}
=\left(\sqrt{\frac{19}{192}}+o(1)\right)
\frac{19^{n/3}}{\sqrt n}$
\end{adjustbox}
\]

as a lower estimate. This improves the existing interval-counting lower estimate by an order-sqrt(n) factor. It does not assert that the sqrt(n) denominator is optimal, and does not change the established exponential rate. The remaining gap to the upper construction is polynomial, not exponential.

\section{2. Exact generating polynomial for distinct ternary sums}\label{n01-exact-generating-polynomial-for-distinct-ternary-sums}

Use the source theorem to choose minimal signed relations r\^{}(j), with disjoint supports B\_j of sizes m\_j \textgreater= 3, one representative from each opposite pair. Let F be the uncovered coordinates. The complete short kernel is precisely

\[
\begin{adjustbox}{max width=.92\linewidth}
$\displaystyle \left\{\sum_j t_jr^{(j)}:t_j\in\{-2,-1,0,1,2\}\right\}.$
\end{adjustbox}
\]

For a block B, write S\_B = sum\_\{i in B\} a\_i and Q\_B = sum\_\{i in B\} a\_i\^{}2. Since its signed relation has coefficients +/-1 on the block, its positive and negative parts have the same sum S\_B/2; in particular S\_B is even. Define

\[
\begin{adjustbox}{max width=.92\linewidth}
$\displaystyle P_B(z)=\prod_{i\in B}(1+z^{a_i}+z^{2a_i})
-z^{S_B/2}\prod_{i\in B}(1+z^{a_i}).$
\end{adjustbox}
\]

\textbf{Proposition (unweighted polynomial factorization).}

\[
\begin{adjustbox}{max width=.92\linewidth}
$\displaystyle \boxed{
\sum_{u\in T(A)}z^u=
\prod_{i\in F}(1+z^{a_i}+z^{2a_i})\prod_jP_{B_j}(z).
}$
\end{adjustbox}
\]

Each factor P\_B has coefficients zero or one. The full product also has coefficients zero or one; this is an identity of unweighted image polynomials, not just an equality of supports or an equality after evaluation at z = 1.

\textbf{Proof.} Reflect a coordinate x\_i to 2-x\_i where r\_i = -1, leaving positive coordinates unchanged. On the reflected cube the equivalence relation is translation by a constant vector. Each class has exactly one representative with minimum coordinate zero, obtained by subtracting the minimum.

The vectors excluded from that representative set have all reflected coordinates in \{1,2\}. Before reflection, a positive coordinate therefore lies in \{1,2\}, and a negative coordinate in \{0,1\}. Their generating polynomial is

\[
\begin{adjustbox}{max width=.92\linewidth}
$\displaystyle z^{\sum_{r_i=1}a_i}\prod_{i\in B}(1+z^{a_i})
=z^{S_B/2}\prod_{i\in B}(1+z^{a_i}).$
\end{adjustbox}
\]

Subtracting it from the full ternary polynomial selects one representative per class, and gives P\_B. The short-kernel theorem ensures that two local representatives have the same value exactly when they are the same class. It also ensures that an equality of global sums cannot change a free coordinate or the class of any block. Thus the Cartesian product of all local representative sets maps bijectively onto T(A). Multiplying their polynomials proves the proposition. QED.

\section{3. Exact variance of the uniform distinct-image distribution}\label{n01-exact-variance-of-the-uniform-distinct-image-distribution}

Let U be uniformly distributed on the finite set T(A), not the ternary cube.

\textbf{Proposition (exact second moment).}

\[
\begin{adjustbox}{max width=.92\linewidth}
$\displaystyle \mathbb E U=S,\qquad
\operatorname{Var}(U)=
\frac23\sum_{i\in F}a_i^2+
\sum_j\kappa_{m_j}Q_{B_j},$
\end{adjustbox}
\]

where

\[
\begin{adjustbox}{max width=.92\linewidth}
$\displaystyle \kappa_m=
\frac{(2/3)3^m-(1/4)2^m}{3^m-2^m}
=\frac{8\,3^m-3\,2^m}{12(3^m-2^m)}.$
\end{adjustbox}
\]

\textbf{Proof.} In a block of size m, the 3\^{}m ternary assignments, counted with multiplicity, have mean S\_B and variance (2/3)Q\_B. The 2\^{}m assignments subtracted in the preceding proposition have the distribution of S\_B/2 plus a sum of independent fair binary choices of the weights a\_i. They have the same mean S\_B and variance Q\_B/4. Subtract the unnormalized zeroth, first, and second moments of these two multisets. The resulting polynomial has \(3^m-2^m\) coefficients equal to one, so division by that number gives mean S\_B and variance kappa\_m Q\_B for the uniform distinct block image.

The global representative map is a bijection. Uniform measure on its domain is therefore the product of the uniform local measures, and pushes forward to uniform measure on T(A). These independent contributions have additive means and variances. A free coordinate has variance 2a\_i\^{}2/3. QED.

For every m \textgreater= 3,

\[
\begin{adjustbox}{max width=.92\linewidth}
$\displaystyle \frac{16}{19}-\kappa_m
=\frac{5(8\,3^m-27\,2^m)}{228(3^m-2^m)}\ge0,$
\end{adjustbox}
\]

because (3/2)\^{}m \textgreater= (3/2)\^{}3. Equality holds at m = 3. Also 2/3 \textless{} 16/19. Consequently

\[
\begin{adjustbox}{max width=.92\linewidth}
$\displaystyle \operatorname{Var}(U)\le\frac{16}{19}Q.$
\end{adjustbox}
\]

The distinction between the two probability spaces matters. For A = \{1,4,5\}, the ternary cube has 27 assignments but only 19 distinct sums. The cube variance is 28, while the uniform distinct-image variance is 672/19. Using 28 for the latter distribution would invalidate the argument.

\section{4. A discrete variance packing inequality}\label{n01-a-discrete-variance-packing-inequality}

\textbf{Lemma.} If V is uniform on M distinct integers, then

\[
\begin{adjustbox}{max width=.92\linewidth}
$\displaystyle \operatorname{Var}(V)\ge\frac{M^2-1}{12}.$
\end{adjustbox}
\]

\textbf{Proof.} List the integers as v\_1 \textless{} ... \textless{} v\_M. Then v\_j-v\_i \textgreater= j-i for j\textgreater i, and

\[
\begin{adjustbox}{max width=.92\linewidth}
$\displaystyle M^2\operatorname{Var}(V)
=\sum_{i<j}(v_j-v_i)^2
\ge\sum_{d=1}^{M-1}(M-d)d^2
=\frac{M^2(M^2-1)}{12}.$
\end{adjustbox}
\]

The last equality follows by the finite formulas for sums of squares and cubes, or directly by induction on M. The M = 1 case is also included. QED.

Applying the lemma to the preceding exact variance yields

\[
\begin{adjustbox}{max width=.92\linewidth}
$\displaystyle \frac{M^2-1}{12}\le\operatorname{Var}(U)\le\frac{16}{19}Q,$
\end{adjustbox}
\]

which proves 19(M\^{}2-1) \textless= 192Q. No density heuristic, concentration approximation, or assumption of distinct ternary representations is used.

\section{5. Exact residue-sensitive optimization of the class count}\label{n01-exact-residue-sensitive-optimization-of-the-class-count}

The source quotient gives M = 3\^{}\textbar F\textbar{} product\_j(\(3^{m_j}-2^{m_j}\)). For m \textgreater= 3,

\[
\begin{adjustbox}{max width=.92\linewidth}
$\displaystyle 3^m-2^m\ge19\,3^{m-3},$
\end{adjustbox}
\]

since this is equivalent to 8*3\^{}\{m-3\} \textgreater= 2\^{}m. The inequality is strict for m \textgreater{} 3. If h is the number of blocks, h \textless= q = floor(n/3), so

\[
\begin{adjustbox}{max width=.92\linewidth}
$\displaystyle M\ge19^h3^{n-3h}\ge19^q3^r=M_n,$
\end{adjustbox}
\]

where the second inequality uses 19 \textless{} 27. Equality M = M\_n holds precisely when there are q blocks, all of size three, and r free coordinates. This characterizes equality in the class-count bound, not in the variance bound or in the extremal problem for N. Combining this with Section 4 proves the theorem in Section 1.

\section{6. Distinctness-aware, integer-only finite bound}\label{n01-distinctness-aware-integer-only-finite-bound}

Distinct positive generators bounded by N satisfy

\[
\begin{adjustbox}{max width=.92\linewidth}
$\displaystyle S\le nN-\frac{n(n-1)}2,\qquad
Q\le Q_{\max}(n,N):=
 nN^2-n(n-1)N+\frac{n(n-1)(2n-1)}6.$
\end{adjustbox}
\]

Thus every admissible N satisfies all three integer constraints

\[
\begin{adjustbox}{max width=.92\linewidth}
$\displaystyle N\ge n,\qquad
2nN\ge M_n+n(n-1)-1,\qquad
192Q_{\max}(n,N)\ge19(M_n^2-1).$
\end{adjustbox}
\]

The companion script defines \texttt{finite\_lower(n)} as the least integer satisfying them. The quadratic expression is increasing for N \textgreater= n; binary search therefore computes the exact threshold without floating-point square roots. For reference, completing the square gives the quadratic threshold

\[
\begin{adjustbox}{max width=.92\linewidth}
$\displaystyle N\ge\frac{n-1}{2}+
\sqrt{\max\left\{0,
\frac{19(M_n^2-1)}{192n}-\frac{n^2-1}{12}\right\}},$
\end{adjustbox}
\]

used together with N \textgreater= n and the interval constraint. When the radicand before taking its maximum is negative, the quadratic adds no restriction.

\begin{longtable}[]{@{}rr@{}}
\toprule\noalign{}
n & Combined finite lower bound \\
\midrule\noalign{}
\endhead
\bottomrule\noalign{}
\endlastfoot
1 & 1 \\
2 & 3 \\
3 & 5 \\
4 & 11 \\
5 & 27 \\
6 & 49 \\
9 & 724 \\
12 & 11,841 \\
30 & 352,128,817,514 \\
\end{longtable}

The first three bounds are attained by \{1\}, \{1,3\}, and \{1,4,5\}, respectively; their binary subset-sum sets can be inspected directly. No claim of exactness is made for the other entries.

\section{7. Reproducibility and formalization boundary}\label{n01-reproducibility-and-formalization-boundary}

Run from the repository root:

\begin{Verbatim}[breaklines=true,breakanywhere=true,fontsize=\footnotesize]
python certificates/verify_variance_lower_bound.py \
  --output certificates/ternary_variance_receipt.json
\end{Verbatim}

The default exhaustive domain is all nonempty subsets of \{1,...,24\} with size at most five: 55,454 candidate sets, of which 2,551 are admissible. For each admissible set the script independently enumerates the entire short kernel, reconstructs the minimal blocks, compares the full polynomial coefficient map against the directly computed unweighted ternary image, checks its exact rational mean and variance, and checks the finite bounds and the class-count equality criterion. This includes 823 cases with a size-three block and 179 with a size-four block; 1,549 have no signed block.

Seven separate targeted cases cover a pair of size-three blocks, a size-three block with two free coordinates, a mixed size-three/size-four pair, and single blocks of sizes three through six. Further bounded checks cover canonical cube moments through dimension eight, the partition optimizer through n=200, and the coefficient bound through m=256. Six regression checks distinguish weighted from unweighted variance, detect an incorrect polynomial shift, and reject inadmissible or invalid inputs. All calculations use integers or \texttt{fractions.Fraction}. The receipt hashes the executable and the ordered exhaustive transcript. Finite checks are corroboration, not the general proof.

No Lean executable was available in the session environment. A new local Lean build was therefore not run, and no uncompiled Lean file is presented as a formal certificate. The existing PR \#1 receipt was inspected, not regenerated.

The next formalization units are precise:

\begin{enumerate}
\def\labelenumi{\arabic{enumi}.}
\tightlist
\item
  Prove the minimal-block and complete short-kernel classification already stated in the source manuscript, including disjoint magnitude layers.
\item
  Construct a finite equivalence from the product of canonical local representatives to the distinct ternary image; this is the essential cardinality and probability-measure bridge.
\item
  Prove the polynomial identity and rational centered-second-moment identity.
\item
  Prove the integer spacing/variance inequality and combine it with the residue-sensitive product bound to obtain the finite lower theorem.
\item
  Add the real nth-root squeeze only after those finite claims are checked.
\end{enumerate}

These are targets, not assertions of completed formalization. The accepted Lean upper-bound theorem, Core hashes, and existing verification statuses remain untouched.

\section{References and dependencies}\label{n01-references-and-dependencies}

Anonymous/deleted Reddit contributor, \& The Clankers. (2026, September 13). \emph{The exponential rate for four-term-progression-free subset-sum sets} {[}Workbench manuscript{]}. \texttt{paper/ep817\_k4\_rate.tex} at the base commit above. The original proof route belongs to the anonymous contributor; reconstruction and verification are attributed separately in that manuscript.

sneed-and-feed. (2026, September 13). \emph{Formalize finite upper bound theorem and core helper lemmas in Extended.lean} {[}Pull request \#1{]}. \nolinkurl{https://github.com/KokunoYumeto/erdos-problem-817-workbench/pull/1}.

Korsky, S. (2026). \emph{Arithmetic progression-free subset-sum sets} {[}Preprint{]}. arXiv:2606.24139. Cited for the surrounding problem and published abstract bounds, not as a source of the new variance formula.
